\scanpage{99}

\mspage{105}

\uline{Définition 3} : On dit que $f:X\to Y$ dans $\mathrm{Diag}$ est une $\mathbb D$-\uline{fibration complète} si c'est une $\mathbb D$-fibration cohomologique \inserted{complète} et une $\mathbb D$-fibration homologique \inserted{complète}.

C'est stable par chgt de base, par composition,

\deleted{\unsure{Comme} \unsure{corol.} : \unsure{les}}

La définition 3 équivaut à ceci : pour tout $p:Y''\to Y'$ sur $Y$, avec $Y''$ et $Y'$ Hot asphériques, le morphisme correspondant $X''\to X'$ est \uline{à la fois} une $\mathbb D$-équivalence cohomologique, et une $\mathbb D$-équivalence homologique. (Et il suffit \unsure{$=$} de le supposer pour $Y''=e$, et quand $Y'$ a un objet final ou un objet initial.) On peut songer à renforcer encore cette condition, en exigeant (sans regarder à gauche ni à droite) que $X''\to X'$ soit une \unsure{$\mathbb D$|$D$}-équivalence $1$-complète.

Mais il n'est pas clair que cette notion renforcée soit stable par composition -- à moins de la renforcer encore, en exigeant que si $Y''\to Y'$ est une $\mathbb D$-équiv. \unsure{$1$-complète|l-complète}, alors $X''\to X'$ l'est aussi. L'ennui avec cette définition renforcée, c'est que faute d'avoir une description cohomologique ou homologique de la \deleted{\unsure{coh.}}notion de

% archive sheet number: 98
